I dette afsnit beskrives dataopdelingen, og hvordan modelarkitekturen vælges.


\subsubsection{Opdeling af data til eksperiment}
Formålet med eksperimenterne er at undersøge modellens evne til at forudsige den næste tidstilstand, iterativt approksimere tidsudviklingsoperatoren,
generalisere til både kendte og ukendte fysiske begyndelsesbetingelser samt overføre viden mellem simuleringer med forskellige rumlige opløsninger.

For at muliggøre disse undersøgelser anvendes kun datasættene \texttt{Alexander\_phi\_0} og \texttt{Alexander\_e} til træning, mens de resterende simuleringer reserveres til test.
Dette gør det muligt at evaluere modellens generaliseringsevne på fysiske scenarier, som ikke er observeret under træningen.

Datasættet \texttt{Alexander\_pi} indeholder en længere tidsserie end både \texttt{Alexander\_phi\_0} og \texttt{Alexander\_e}. Det anvendes derfor som et
særligt testsæt til at undersøge modellens evne til at forudsige plasmaets udvikling over længere tidshorisonter end dem, der indgår i træningsdataene.

Da modellen både skal kunne interpolere og ekstrapolere i tidsdimensionen, opdeles datasættene forskelligt. Tidsrækken i \texttt{Alexander\_e}
opdeles kronologisk, således at de tidligste observationer anvendes til træning, mens senere observationer anvendes til validering og test.
For \texttt{Alexander\_phi\_0} foretages derimod en tilfældig opdeling af samples.

For begge datasæt anvendes 70\% af data til træning, 15\% til validering og 15\% til test.

\subsubsection{Domæne for modelarkitektur}
Domænet for antal lag, aktiveringsfunktioner og hidden dimension blev diskuteret i afsnit \ref{ch2_4_2}. På den baggrund undersøges følgende kombinationer:
\begin{equation*}
\begin{aligned}
    \text{Antal lag} &\in \{2, 4, 6, 8\}, \\
    \text{Aktiveringsfunktion} &\in \{\tanh, \text{SiLU}, \text{blanding af tanh og SiLU}\}, \\
    \text{Hidden dimension} &\in \{256, 512, 1024, 2048, 4096\}.
\end{aligned}
\end{equation*}
Først kommer Hidden dimension til at være låst fast til 256. Så det kun er antallet af lag og aktiveringsfunktion der varierer.
Herefter vælges nogle af de bedste aktiveringsfunktion og antal lag, og så undersøges det bedste hidden dimension konfiguration.

Da SiLU aktiveringsfunktionen er yderst sensitiv overfor learning rate, kommer learning rate til at variere imellem eksperimenterne. Når der kommer
flere lag i netværket, bliver blandingen af aktiveringsfunktioner opdelt i fire grupper, første gruppe har tanh og SiLU skiftevis, 
anden gruppe har SiLU og tanh skiftevis tedje gruppe har første halvdel tanh og anden halvdel SiLU, fjerde gruppe har første halvdel SiLU og anden halvdel tanh.
Dette gøres for at give begge aktiveringsfunktioner mulighed for at konvergere.
